\begin{onehalfspace}

Il capitolo precedente ha descritto l'architettura del sistema e la logica dei due metodi. Prima di entrare nell'implementazione, è necessario fissare con precisione cosa il sistema deve fare, quali grandezze prende in ingresso, quali produce in uscita e come si giudica se un risultato è valido. Questo capitolo ha questo scopo: fornisce le specifiche funzionali del sistema sperimentale, definisce i quattro use case su cui vengono condotti gli esperimenti e stabilisce le metriche con cui vengono valutati e confrontati i due metodi.

% -----------------------------------------------
\section{Requisiti Funzionali}
\label{sec:requisiti}

Il sistema sperimentale deve soddisfare i seguenti requisiti funzionali:

\begin{enumerate}

    \item \textbf{Esecuzione configurabile del circuito di Shor.}
    Dato un intero $N$ da fattorizzare e una base $a$ con $\gcd(a, N) = 1$, il sistema costruisce il circuito dell'algoritmo di Shor a partire dall'implementazione di riferimento \cite{github_shor}, lo configura con il numero di qubit $n_{\text{count}}$ specificato e lo esegue tramite \texttt{AerSimulator} con il noise model attivo per un numero di shot $M_{\text{shots}}$ configurabile.

    \item \textbf{Modello di rumore parametrizzabile.}
    Il sistema permette di configurare in modo indipendente tutti i parametri del noise model --- tassi di errore depolarizzante per gate a singolo e doppio qubit, tempi di rilassamento $T_1$ e $T_2$, probabilità di readout error --- in modo da poter riprodurre diverse condizioni operative e generare dataset con distribuzione di rumore variabile. Questo requisito è fondamentale per la creazione di un dataset di addestramento sufficientemente ricco per il Metodo 2, come indicato dalla metodologia di selezione dei noise model proposta in letteratura \cite{epjqt_noise2024}.

    \item \textbf{Raccolta e archiviazione degli output.}
    Il sistema registra la distribuzione di misura del registro di controllo per ogni esecuzione e la archivia in formato strutturato, garantendo la riproducibilità degli esperimenti tramite seed fisso per il generatore di numeri casuali di Qiskit Aer.

    \item \textbf{Analisi statistica del Metodo 1.}
    Il sistema implementa la pipeline del Metodo 1: identificazione dei picchi della distribuzione tramite filtraggio a soglia, stima del periodo $r$ tramite frazioni continue, verifica della correttezza dei fattori estratti. Misura e registra il numero di iterazioni necessarie per raggiungere il primo risultato corretto.

    \item \textbf{Training, validazione e inferenza del Metodo 2.}
    Il sistema genera il dataset dal simulatore rumoroso, addestra il classificatore con divisione training/validation/test, ne valuta le metriche di accuratezza e lo utilizza in modalità inferenza: dopo ogni esecuzione del circuito, il classificatore decide se il risultato è corretto e, in caso affermativo, interrompe il processo.

    \item \textbf{Confronto sistematico sui quattro use case.}
    Entrambi i metodi vengono eseguiti sugli stessi quattro use case, nelle medesime condizioni di rumore, producendo metriche confrontabili secondo i criteri definiti nella Sezione~\ref{sec:metriche}.

    \item \textbf{Riproducibilità.}
    Tutti gli esperimenti sono riproducibili: seed fisso, parametri esplicitati, codice disponibile nel repository del progetto. Questo requisito è condizione necessaria per la validazione accademica dei risultati \cite{shor_hpec2023}.

\end{enumerate}

% -----------------------------------------------
\section{Parametri di Input e Output}
\label{sec:input_output}

\subsection{Parametri di Input}

La Tabella~\ref{tab:input_params} elenca i parametri di configurazione del sistema. I loro valori vengono fissati per ciascun use case prima dell'esecuzione e rimangono costanti per tutta la durata dell'esperimento corrispondente.

\begin{table}[H]
\centering
\begin{tabular}{@{}lll@{}}
\toprule
\textbf{Parametro}  & \textbf{Unità} & \textbf{Descrizione} \\
\midrule
$N$                 & intero         & Numero da fattorizzare \\
$a$                 & intero         & Base della moltiplicazione modulare ($\gcd(a,N)=1$) \\
$n_{\text{count}}$  & intero         & Qubit del registro di controllo \\
$M_{\text{shots}}$  & intero         & Shot per esecuzione del simulatore \\
$\epsilon_{1q}$     & adimensionale  & Tasso di errore depolarizzante su gate a 1 qubit \\
$\epsilon_{2q}$     & adimensionale  & Tasso di errore depolarizzante su gate a 2 qubit \\
$T_1$               & ns             & Tempo di rilassamento (amplitude damping) \\
$T_2$               & ns             & Tempo di dephasing \\
$p_{\text{ro}}$     & adimensionale  & Probabilità di readout error \\
\bottomrule
\end{tabular}
\caption{Parametri di input del sistema sperimentale. I valori di $\epsilon_{1q}$, $\epsilon_{2q}$, $T_1$, $T_2$ e $p_{\text{ro}}$ definiscono il noise model e vengono variati tra i use case per coprire scenari di rumore diversi.}
\label{tab:input_params}
\end{table}

\subsection{Output del Sistema}

Per ogni use case e per ciascuno dei due metodi, il sistema produce le seguenti quantità:

\begin{itemize}
    \item I fattori trovati $p$ e $q$ tali che $p \cdot q = N$, con verifica esplicita della correttezza ($p \neq 1$, $q \neq 1$, $p \cdot q = N$);
    \item Il numero di iterazioni necessarie per il primo risultato corretto;
    \item La distribuzione di frequenza completa degli output del registro di controllo su $M_{\text{shots}}$ esecuzioni;
    \item Le metriche di confronto tra i due metodi definite nella Sezione~\ref{sec:metriche}.
\end{itemize}

Per il solo Metodo 2, vengono prodotte in aggiunta:
\begin{itemize}
    \item L'accuratezza del classificatore e il punteggio F1 sul test set;
    \item La curva ROC con l'area sottostante (\ac{AUC}).
\end{itemize}

% -----------------------------------------------
\section{I Quattro Use Case}
\label{sec:use_case}

% TODO [AGGIORNARE DOPO I PRIMI ESPERIMENTI]
% I parametri di rumore e i valori di n_count riportati in questa sezione sono
% stati fissati a partire dalla letteratura e dalle specifiche di ibm_marrakesh
% PRIMA di eseguire gli esperimenti. Una volta ottenuti i primi risultati sul
% simulatore, verificare che:
%   (a) i circuiti per N=21 e N=35 vengano costruiti correttamente dal repo;
%   (b) i livelli di rumore producano tassi di successo nell'intervallo atteso
%       (5-40% per il livello realistico, <5% per il livello degradato);
%   (c) n_count sia sufficiente a risolvere il periodo r per ciascun caso.
% Se necessario, aggiornare la tabella Tab. \ref{tab:use_case_params} con i
% valori effettivamente usati durante le esecuzioni.

I quattro use case sono scelti per garantire che i risultati non dipendano da una singola configurazione e che coprano un intervallo rappresentativo di condizioni operative. Il disegno sperimentale segue due assi di variazione indipendenti:

\begin{itemize}
    \item \textbf{Asse rumore (Use Case 1 vs.\ 2)}: stesso numero $N = 15$, stessa base $a = 7$, stesso circuito --- livello di rumore diverso. Il confronto isola l'effetto del rumore sul numero di iterazioni necessarie, a parità di complessità circuitale. Il Use Case 2 utilizza parametri di errore circa cinque volte superiori a quelli del Use Case 1, per coprire sia un regime \textit{NISQ-realistico} sia uno \textit{NISQ-degradato}.

    \item \textbf{Asse scalabilità (Use Case 1, 3, 4)}: numeri $N$ di dimensione crescente ($15 \to 21 \to 35$), stesso livello di rumore NISQ-realistico. Il confronto valuta come il numero medio di iterazioni e l'accuratezza del classificatore varino al crescere della profondità del circuito.
\end{itemize}

La scelta di quattro use case soddisfa il requisito metodologico minimo per la validazione accademica dei risultati in questo dominio \cite{shor_hpec2023}.

\subsection{Parametri dei Use Case}

La Tabella~\ref{tab:use_case_params} riporta la configurazione completa dei quattro use case. I parametri di rumore del livello \textit{NISQ-realistico} sono derivati dalle specifiche di calibrazione di \texttt{ibm\_marrakesh} riportate nel tutorial ufficiale IBM \cite{ibm_shor_tutorial} e dalla metodologia di selezione dei noise model proposta in \cite{epjqt_noise2024}. Il livello \textit{NISQ-degradato} scala i tassi di errore di un fattore $5\times$ e dimezza i tempi di coerenza, simulando condizioni di hardware più antico o in stato di calibrazione peggiore.

\begin{table}[H]
\centering
\resizebox{\textwidth}{!}{%
\begin{tabular}{@{}lcccc@{}}
\toprule
\textbf{Parametro} & \textbf{UC 1} & \textbf{UC 2} & \textbf{UC 3} & \textbf{UC 4} \\
\midrule
$N$                          & 15  & 15  & 21  & 35  \\
$a$                          & 7   & 7   & 2   & 6   \\
Periodo $r$ atteso           & 4   & 4   & 6   & 2   \\
$n_{\text{count}}$ (qubit)   & 8   & 8   & 10  & 12  \\
$M_{\text{shots}}$           & 4096 & 4096 & 4096 & 4096 \\
\midrule
\textit{Livello rumore}      & NISQ-realistico & NISQ-degradato & NISQ-realistico & NISQ-realistico \\
$\epsilon_{1q}$              & $1\times10^{-3}$ & $5\times10^{-3}$ & $1\times10^{-3}$ & $1\times10^{-3}$ \\
$\epsilon_{2q}$              & $1\times10^{-2}$ & $5\times10^{-2}$ & $1\times10^{-2}$ & $1\times10^{-2}$ \\
$T_1$ (ns)                   & 100\,000 & 50\,000 & 100\,000 & 100\,000 \\
$T_2$ (ns)                   & 80\,000  & 30\,000 & 80\,000  & 80\,000  \\
$p_{\text{ro}}$              & $2\times10^{-2}$ & $5\times10^{-2}$ & $2\times10^{-2}$ & $2\times10^{-2}$ \\
\bottomrule
\end{tabular}%
}
\caption{Configurazione completa dei quattro use case. I parametri di rumore del livello NISQ-realistico sono derivati dalle specifiche di \texttt{ibm\_marrakesh} \cite{ibm_shor_tutorial, epjqt_noise2024}. Il livello NISQ-degradato scala i tassi di errore di $5\times$ e dimezza $T_1$ e $T_2$.}
\label{tab:use_case_params}
\end{table}

\subsection{Razionale per la Scelta delle Istanze}

\textbf{Use Case 1 --- $N=15$, NISQ-realistico.}
È il caso di riferimento consolidato in letteratura: la prima dimostrazione sperimentale di Shor è avvenuta su $N=15$ con NMR a 7 qubit \cite{vandersypen2001}, e il tutorial ufficiale IBM riporta i dati di esecuzione per la stessa istanza su \texttt{ibm\_marrakesh} \cite{ibm_shor_tutorial}. La base $a=7$ produce un periodo $r=4$, sufficiente per estrarre i fattori $3$ e $5$ tramite l'algoritmo delle frazioni continue. Questo use case funge da baseline quantitativa per entrambi i metodi.

\textbf{Use Case 2 --- $N=15$, NISQ-degradato.}
Configurazione identica al Use Case 1 eccetto per il livello di rumore, aumentato di $5\times$ su tutti i parametri. Poiché il circuito è lo stesso, qualsiasi differenza nelle metriche $\bar{M}_1$ e $\bar{M}_2$ è attribuibile esclusivamente al rumore. Questo use case permette di valutare la \textbf{robustezza} del classificatore ML rispetto al degrado delle condizioni hardware, indipendentemente dalla complessità circuitale.

\textbf{Use Case 3 --- $N=21$, NISQ-realistico.}
$N = 21 = 3 \times 7$ è la fattorizzazione più grande dimostrata su hardware superconduttore a uso generale \cite{skosana2021}. La base $a=2$ produce un periodo $r=6$, che implica una profondità circuitale sensibilmente maggiore rispetto a $N=15$. A parità di livello di rumore con il Use Case 1, questo use case isola l'effetto della scalabilità: ci si attende un aumento di $\bar{M}_1$ per via del circuito più profondo, e si valuta se il classificatore ML mantiene un'accuratezza sufficiente nonostante la maggiore complessità.

\textbf{Use Case 4 --- $N=35$, NISQ-realistico.}
$N = 35 = 5 \times 7$ con base $a=6$ produce il periodo minimo $r=2$ (poiché $6^2 = 36 \equiv 1 \pmod{35}$), il che consente di fattorizzare direttamente con $\gcd(6-1,\, 35) = \gcd(5, 35) = 5$ e $\gcd(6+1,\, 35) = \gcd(7, 35) = 7$. La scelta di $r=2$ con $N$ maggiore permette di separare l'effetto della dimensione del registro (più qubit) dall'effetto della profondità della modular exponentiation, fornendo un punto di confronto complementare rispetto al Use Case 3.

% -----------------------------------------------
\section{Metriche di Valutazione}
\label{sec:metriche}

\subsection{Metriche del Metodo 1}

\begin{itemize}
    \item $\bar{M}_1$: numero medio di esecuzioni necessarie per ottenere il primo risultato corretto, calcolato su $K = 50$ ripetizioni indipendenti dello stesso esperimento.
    \item $\sigma_{M_1}$: deviazione standard del numero di esecuzioni, a indicare la variabilità del metodo.
    \item $P_{\text{success},1}(M_{\text{max}})$: probabilità di successo entro un budget fisso di $M_{\text{max}}$ esecuzioni, dove $M_{\text{max}}$ verrà fissato in base ai primi risultati sperimentali.
\end{itemize}

\subsection{Metriche del Metodo 2}

\begin{itemize}
    \item $\bar{M}_2$: numero medio di esecuzioni prima che il classificatore restituisca la prima etichetta $1$ corretta, sulle stesse $K$ ripetizioni usate per il Metodo 1.
    \item $\sigma_{M_2}$: deviazione standard corrispondente.
    \item $\text{Acc}_{\text{clf}}$: accuratezza del classificatore sul test set, definita come il rapporto tra predizioni corrette e numero totale di campioni di test.
    \item $\text{F1}_{\text{clf}}$: F1-score, che bilancia la precisione e il recall nella classe positiva (output corretto). È la metrica preferita rispetto alla sola accuratezza in presenza di classi sbilanciate, un fenomeno frequente quando il tasso di successo dell'algoritmo è basso.
    \item $P_{\text{success},2}(M_{\text{max}})$: probabilità di successo entro lo stesso budget $M_{\text{max}}$ utilizzato per il Metodo 1.
\end{itemize}

\subsection{Metrica di Confronto}

La metrica principale per il confronto tra i due metodi è il \textbf{fattore di riduzione delle iterazioni}:
\begin{equation}
    \rho = \frac{\bar{M}_1}{\bar{M}_2}
    \label{eq:fattore_riduzione}
\end{equation}
Un valore $\rho > 1$ indica che il Metodo 2 converge più rapidamente del Metodo 1. L'ipotesi sperimentale, motivata nel Capitolo~\ref{chap:strategie}, è $\rho \gg 1$, con $\bar{M}_2 \approx 3$. La significatività statistica del risultato verrà valutata tramite un test $t$ di Welch sulle distribuzioni di $M_1$ e $M_2$.

\subsection{Parametri di Error Mitigation Applicati}

Per entrambi i metodi, durante la simulazione vengono applicati i seguenti strumenti di error mitigation:

\begin{itemize}
    \item \textbf{Readout mitigation}: correzione delle probabilità di misura tramite matrice di calibrazione stimata su circuiti dedicati. Riduce la componente sistematica degli errori di lettura senza aggiungere overhead circuitale.
    \item \textbf{Filtraggio a soglia} (Metodo 1): i valori con frequenza inferiore a una percentuale $\tau$ del massimo osservato vengono scartati prima dell'analisi dei picchi. Il valore ottimale di $\tau$ viene selezionato sperimentalmente per ciascun use case.
\end{itemize}

L'effetto di ciascun parametro di mitigazione viene riportato nei capitoli dei risultati, confrontando le metriche ottenute con e senza la mitigazione attiva.

\end{onehalfspace}
